% Copyright 2026 Yoshida Shin
%
% This file is part of the ``Repeating Nines''.
%
% Permission is granted to copy, distribute, and/or modify this document
% under the terms of the GNU Free Documentation License, Version 1.3 or any later version published by the Free Software Foundation;
% with no Front-Cover Texts, no Back-Cover Texts, and one Invariant Section: ``This document was written by Shin Yoshida with the aim of contributing to a fairer and safer world.''
% A copy of the license is included in the section entitled ``GNU Free Documentation License.''

\begin{frame}[c]{}{}
  {\large Proof 0.9999... = 1}
\end{frame}

\begin{frame}[c]{}{}
  Before the proof, please accept the fact that $1 \div 3 \times 3 = 1$.
\end{frame}

\begin{frame}[c]{}{}
  It is obvious to define ``division'' and ``multiplication''.

  \vspace{2ex}

  But I have to explain many things to do that, for example, ``What is number?''.

  \vspace{2ex}

  It is too much here.
\end{frame}

\begin{frame}[c]{}{}
  The definition of $a_n$ includes multiplication.

  \vspace{4ex}

  I cannot calculate lim $a_n$ without multiplication!
\end{frame}

\begin{frame}[c]{}{}
  Let's prove 0.9999... = 1 using our definition.

  \vspace{2ex}

  i.e. prove that x fulfills the conditions if and only if x = 1.

  \vspace{2ex}

  \begin{enumerate}
  \item $\forall n \in \mathbb{N}, a_n < x$
  \item $\forall k < x, \exists n$ $k \leqq a_n$
  \end{enumerate}
\end{frame}

\begin{frame}[c]{}{}
  We need consider cases based on the value of x.

  \vspace{2ex}

  \begin{itemize}
  \item $x < 1$
  \item $x = 1$
  \item $x > 1$
  \end{itemize}
\end{frame}

\begin{frame}[c]{$x < 1$}{}
  x does not fulfill the condition 1

  if $x < 1$

  \vspace{2ex}

  1. $\forall n \in \mathbb{N}, a_n < x$

  \vspace{2ex}

  $\because if$ $n > -\log_{10} (1 - x), x < a_n$
\end{frame}

\begin{frame}[c]{x = 1}{}
  From the definition,

  $a_n = 1 - (0.1)^n$

  \vspace{2ex}

  $\therefore \forall n \in \mathbb{N}, a_n < 1$

  ($\because 0 < (0.1)^n$)
\end{frame}

\begin{frame}[c]{x = 1}{}
  Also,

  \vspace{2ex}

  $\forall k < 1,$

  if $n > -\log_{10} (1 - k), k \leqq a_n$

  \vspace{2ex}

  $\therefore \forall k < 1, \exists n$ $k \leqq a_n$
\end{frame}

\begin{frame}[c]{x = 1}{}
  Therefore, x = 1 fulfills both the following conditions.

  \vspace{2ex}

  \begin{enumerate}
  \item $\forall n \in \mathbb{N}, a_n < x$
  \item $\forall k < x, \exists n$ $k \leqq a_n$
  \end{enumerate}
\end{frame}

\begin{frame}[c]{$x > 1$}{}
  x does not fulfill the condition 2

  if $x > 1$,

  2. $\forall k < x, \exists n$ $k \leqq a_n$

  \vspace{2ex}

  $\because if$ $k = \frac{x + 1}{2}$
  
  $k < x \land \forall n \in \mathbb{N}, a_n < k$

  ($\because a_n < 1 \land 1 < k$)
\end{frame}

\begin{frame}[c]{}{}
  Let's pause and summarize where we are.

  \vspace{2ex}

  \begin{itemize}
    \item First, we introduced a frivolous game to define the limit of $a_n$
    \item Then, we defined the limit of $a_n$ as the winning value in the game
    \item Next, we rewrote the winning strategy using 2 conditions
  \end{itemize}

\end{frame}

\begin{frame}[c]{}{}
  Here, we have proved

  \vspace{4ex}

  ``x = 1 fulfills those 2 conditions.''
\end{frame}

\begin{frame}[c]{}{}
  At the same time, we proved

  \vspace{4ex}

  ``Only if x = 1, x fulfills those conditions.''
\end{frame}

\begin{frame}[c]{}{}
  i.e. we have proved

  \vspace{4ex}

  ``0.9999... = 1.''
\end{frame}
